\documentclass[../../main.tex]{subfiles}
\begin{document}

\begin{center}
  {\fontsize{48pt}{58pt}\selectfont \textbf{2001}\normalsize\textbf{年度}}
\end{center}

\begin{table}[htb]
  \centering
  \caption{出題テーマと難易度}
  \begin{tabular}{|c|p{9cm}|c|}
    \hline
    \textbf{問題番号} & \textbf{テーマ} & \textbf{難易度} \\
    \hline
    第1問 & 正整数解をもつ三変数二次方程式 & \\
    \hline
    第2問 & 多項式の漸化式と収束範囲 & \\
    \hline
    第3問 & 正五角形の延長線と複素数 & \\
    \hline
    第4問 & 常用対数と2の累乗の先頭数字 & \\
    \hline
    第5問 & 連立一次方程式の解の有無 & \\
    \hline
    第6問 & 円の接線から描く曲線の面積 & \\
    \hline
  \end{tabular}
\end{table}

\begin{table}[htb]
  \centering
  \caption{解答}
  \renewcommand{\arraystretch}{1.4}
  \begin{tabular}{|c|c|p{9cm}|}
    \hline
    \textbf{問題番号} & \textbf{小問} & \textbf{答え} \\
    \hline
    第1問 & -- & $(x,y,z)=(3,2,1),(3,1,2)$ \\
    \hline
    第2問 & -- & $g_n(x)=1-x^{2^{n-1}}$；収束範囲$-1\le x<1$ \\
    \hline
    第3問 & (1) & $z=-\alpha^3+\alpha+1$ \\
    \cline{2-3}
    & (2) & 証明問題 \\
    \hline
    第4問 & (1) & $n=10$の一例 \\
    \cline{2-3}
    & (2) & $n=56$の一例 \\
    \hline
    第5問 & (1)，(2) &  \\
    \hline
    第6問 & -- & $C_2$の最小$x$座標は$-1$未満；囲む面積$\dfrac{\pi}{2}(a^2+2)$ \\
    \hline
  \end{tabular}
\end{table}

\paragraph{試験形態}
6問．試験時間は要確認．

\end{document}
